% !TeX root = ../traffic_rules.tex

\section{Preliminaries}
\label{sec:preliminaries}
\subsection{Road Network}
We describe our road network with lanelets, which are atomic, drivable road segments, geometrically represented by a left and right bound, where the bounds are represented by polylines \cite{Bender2014}. The set of all lanelets is denoted by $\mathbb{L}$. In the following we refer to a single lanelet as $l$, where $l \in \mathbb{L}$. 
The functions $\mathrm{l}_\mathrm{lb}(l)$, $\mathrm{l}_\mathrm{rb}(l)$, $\mathrm{l}_\mathrm{begin}(l)$, and $\mathrm{l}_\mathrm{end}(l)$ return the left bound, right bound, left and right beginning points, and left and right ending points of a lanelet. The functions
\begin{equation}
\mathrm{adj}_\mathrm{l}(l) =
\begin{cases} 
l' & \text{if } \exists \ l' \in \mathbb{L}: \mathrm{l}_\mathrm{lb}(l') = \mathrm{l}_\mathrm{rb}(l) \\ 
\emptyset & \text{otherwise} 
\end{cases}	
\end{equation}

\begin{equation}
\mathrm{adj}_\mathrm{r}(l) =
\begin{cases} 
l' & \text{if } \exists \ l' \in \mathbb{L}: \mathrm{l}_\mathrm{rb}(l) = \mathrm{l}_\mathrm{lb}(l') \\ 
\emptyset & \text{otherwise} 
\end{cases}	
\end{equation}

\begin{equation}
\mathrm{succ}(l) =
\begin{cases} 
l' & \text{if } \exists \ l' \in \mathbb{L}: \mathrm{l}_\mathrm{end}(l) = \mathrm{l}_\mathrm{begin}(l') \\ 
\emptyset & \text{otherwise} 
\end{cases}	
\end{equation}

\begin{equation}
\mathrm{pre}(l) =
\begin{cases} 
l' & \text{if } \exists \ l' \in \mathbb{L}: \mathrm{l}_\mathrm{begin}(l) = \mathrm{l}_\mathrm{end}(l') \\ 
\emptyset & \text{otherwise} 
\end{cases}	
\end{equation}
return the adjacent left lanelet, adjacent right lanelet, successor lanelets, and predecessor lanelets of $l$, respectively. The curvature of a lanelet at a given path length $s_\mathrm{l}$ is given by $\mathrm{curv}(s_\mathrm{l})$.The function $\mathtt{speed\_limit}(l)$ returns the maximum speed which a vehicle is allowed to drive on a lanelet. A lanelet can have several attributes which further describe it. The set of all attributes is $\mathbb{A}$.  For the highway case, we have $\mathbb{A} =  \{ main\_carriage\_way, access\_ramp, exit\_ramp, shoulder \}$. The function $attr(l)$ returns the attributes which are valid for a lanelet. Similar, the left and right bounds of a lanelet can have lanelet markings, with the set of lanelet markings $\mathbb{LM} = \{solid, dashed, broad\_solid, broad\_dashed \}$. The functions $\mathrm{mark}_l(l)$ and $\mathrm{mark}_r(l)$ return the left and right lanelet marking, respectively.
We define the lane part in front of a lanelet, in the rear of a lanelet as follows:


\begin{align}
\begin{split}
\mathrm{lane}_\mathrm{front}(l) =  \big\{l' \, \big| \, \forall l'', l''' \in (l' \in (\mathrm{succ}(l) \cup \mathrm{lane}_\mathrm{front}(\mathrm{succ}(l)))): \mathrm{attr}(l'') = \mathrm{attr}(l''')  \big\},
\end{split}
\end{align}

\begin{align}
\begin{split}
\mathrm{lane}_\mathrm{rear}(l) =  \big\{l' \, \big| \, \forall l'', l''' \in (l' \in (\mathrm{pre}(l) \cup \mathrm{lane}_\mathrm{rear}(\mathrm{pre}(l)))): \mathrm{attr}(l'') = \mathrm{attr}(l''')  \big\}.
\end{split}
\end{align}

We define the lane to which a lanelet belongs to, using the lanelets in front and rear of another lanelet:
\begin{align}
\begin{split}
\mathrm{lane}(l) &=  %{l' \cup \mathrm{pre}(l') \cup \mathrm{succ}(l')}
\big\{l' \in (l \ \cup \ \mathrm{lane}_\mathrm{rear}(l)  \ \cup \ \mathrm{lane}_\mathrm{front}(l))   \, \big| \, \forall l'', l''' \in l': \mathrm{attr}(l'') = \mathrm{attr}(l''') \\
& \land \mathrm{adj}_\mathrm{l}(l') \in \mathrm{succ}((\mathrm{adj}_\mathrm{l}(\mathrm{pre}(l'))) \land \mathrm{adj}_\mathrm{r}(l') \in \mathrm{succ}((\mathrm{adj}_\mathrm{r}(\mathrm{pre}(l'))) \big\}.
\end{split}
\end{align}

The left and right lane are defined as follows:
\begin{align}
\mathrm{lane}_\mathrm{left}(l) =\big\{l' \, \big| \, \exists l'' \in \mathrm{lane}(l) : l' = \mathrm{adj}_\mathrm{l}(l'') \big\},
\end{align}
\begin{align}
\mathrm{lane}_\mathrm{right}(l) =\big\{l' \, \big| \, \exists l'' \in \mathrm{lane}(l) : l' = \mathrm{adj}_\mathrm{r}(l'') \big\}.
\end{align}
The functions
\begin{align}
\mathrm{lanes}_\mathrm{left}(l) =\big\{l' \, \big| \, \exists l'' \in \mathrm{lane}(l) : l' = \mathrm{adj}_\mathrm{l}(l'') \lor \exists l''' \in \mathrm{lanes}_\mathrm{left}(l'')  : l' = \mathrm{adj}_\mathrm{l}(l''') \big\},
\end{align}
\begin{align}
\mathrm{lanes}_\mathrm{right}(l) =\big\{l' \, \big| \, \exists l'' \in \mathrm{lane}(l) : l' = \mathrm{adj}_\mathrm{r}(l'') \lor \exists l''' \in \mathrm{lanes}_\mathrm{right}(l'')  : l' =\mathrm{adj}_\mathrm{r}(l''') \big\}
\end{align}
return all left and right lanes of a given lanelet.

\subsection{Vehicle}
A vehicle $V_i$, $i \in \mathbb{N}^+$, consists of a longitudinal, lateral, and signal state $x_\mathrm{lon} \in \mathbb{R}^n$ , $x_\mathrm{lat} \in \mathbb{R}^n$, and $x_\mathrm{sig} \subseteq \mathbb{SIG}$, respectively, where $\mathbb{SIG} = \{horn, indicator\_left,  indicator\_right, braking\_lights, hazard\_warning\_lights, \\ flashing\_blue\_lights\}$ represents possible active signals of a vehicle. The set of all vehicles is denoted by $\mathbb{V}$.The longitudinal state $x_\mathrm{lon} = \begin{bmatrix} s & v & a & j\end{bmatrix}^T$ consists of position $s$, velocity $v$, acceleration $a$, and jerk $j$. 
 %The longitudinal state of other vehicles $x_\mathrm{o} = \begin{bmatrix} s & v & a\end{bmatrix}^T$ consists only of position, velocity , acceleration. 
The lateral state $x_\mathrm{lat} = \begin{bmatrix} d & \theta & \kappa & \dot{\kappa} \end{bmatrix}^T$ consists of the lateral distance to the reference path $\Gamma$, orientation $\theta$, curvature $\kappa$, and change of curvature $\dot{\kappa}$. The longitudinal input is $u(t) = \ddot{a}(t)$ and the lateral input is $u(t) = \ddot{\kappa}(t)$. The signal state contains only the subset of currently active elements of $\mathbb{SIG}$. The occupancy of a vehicle for given states $x_\mathrm{lon}$ and $x_\mathrm{lat}$ is denoted by $\mathcal{O}(x_\mathrm{lon}(t), x_\mathrm{lat}(t))$. The function $\mathrm{proj()}$ projects a vehicle to a position $s$. The functions  $\mathrm{front()}$ and $\mathrm{rear()}$ return the front and rear position of a vehicle. All variables associated to the ego vehicle are denoted by the subscript $\square_\mathrm{ego}$ and all variables associated to other vehicles are denote by the subscript $\square_\mathrm{o}$.
For a given initial state, an input trajectory $u(\cdot)$, we introduce the solution of the model $\dot{x}=f(x,u)$ of the ego vehicle over time $t$ as $\xi(t;x_0,u(\cdot))$. 
We denote a braking input by $u_{brake}(\cdot)$.  The time at which a safe state is reached for $u(\cdot)$, given the system dynamics and $x_0$, is denoted by $t_s(u(\cdot))$; we consider a state to be safe when one can stay in it indefinitely without causing a collision \cite{Pek2018b}.

%
%\begin{definition}[Safe distance]
%\label{def:safe_distance}
%The ACC vehicle can definitely avoid a rear-end collision, if it maintains at least the minimal safe distance $\safeDistance$ to the leading vehicle: %, \ie $\distance(t_0) \geq \safeDistance(t_0)$:
%\begin{align*}
%\safeDistance(t_k) := \inf \big( \big\{ \distance(t_k) \, \big| \, 
%&\forall t_i \in \{t_{k}, t_{k+1}, \ldots, \tegoStop\} : \distance(t_i) > 0 \big\} \big),
%\end{align*}
%where $\distance(t_i)$ is obtained by simulating the leading and ACC vehicle according to the vehicle model from $v(t_k)$ with $u_\lead^\text{brake}\big([t_k, \tegoStop]\big)$ and $u_\ego^\text{brake}\big([t_k, \tegoStop]\big)$, respectively.
%\end{definition}



\subsection{Linear Temporal Logic}
%Given is a set $\Pi$ of atomic propositions, where each atomic proposition $\pi_i \in \Pi$ represents a Boolean statement, and a LTL formula $\phi$. The Backus-Naur form of a LTL formula is
%
%\begin{equation*}
%	\phi ::= \pi_i \,|\, \neg \phi \,|\, \phi_1 \vee \phi_2 \,|\, X(\phi) \,|\, \phi_1 U \phi_2 
%\end{equation*}
%
%\textbf{LTL operators}: The letter X stands for \textit{ne\textbf{x}t state} and U for \textit{\textbf{u}ntil}. Further formulas are later derived from those.
%
%\textbf{Precedence}: unary connectives have precedence over binary connectives, $\neg$ and $X()$ have equal precedence and $U$ has precedence over $\vee$.
%
%The semantics of LTL is defined in terms of a infinite sequence $\sigma = z(0), z(1), z(2), \ldots$, where each $z(i) \in 2^\Pi$:
%
%\begin{tabular}{lcl}
%$\mathcal{K},\pi \vDash p$ & iff & $p \in L(z(0))$ (``The model $\mathcal{K}$ fulfills $p$ for the path $\pi$'').  \\
% $\mathcal{K},\pi \vDash \neg \phi$ & iff & not $\mathcal{K},\pi \vDash \phi$. \\
% $\mathcal{K},\pi \vDash \phi_1 \vee \phi_2$ & iff & $\mathcal{K},\pi \vDash \phi_1$ or $\mathcal{K},\pi \vDash \phi_2$. \\
% $\mathcal{K},\pi \vDash X(\phi)$ & iff & $\mathcal{K},\pi^1 \vDash \phi$.  \\
% $\mathcal{K},\pi \vDash \phi_1 U \phi_2$ & iff & $\exists i\geq 0 \text{ such that } \mathcal{K},\pi^i \vDash \phi_2 \text{ and } \forall j, 0\leq j< i, \mathcal{K},\pi^j \vDash \phi_1$. 
%\end{tabular}\\
%\vspace{0.1cm} \mbox{}
%(iff: if and only if; $\pi^i= z(i), z(i+1), z(i+2), \ldots$)
%
%\textbf{Abbreviations:} \\
%\vspace{0.2cm}
%\begin{tabular}{lcll}
% $true$ & $=$ & $\pi_i \vee \neg \pi_i$  \\
% $false$ & $=$ & $\neg true$  \\
% $\phi_1 \wedge \phi_2$ & $=$ & $\neg (\neg\phi_1 \vee \neg\phi_2)$  \\
% $\phi_1 \implies \phi_2$ & $=$ & $\neg \phi_1 \vee \phi_2$ \\[0.2cm]
% $\phi_1$ R $\phi_2$ & $=$ & $\neg(\neg \phi_1 U \neg \phi_2)$  & \hspace{-0.4cm} $\phi_2$ is true until it is \textbf{r}eleased by $\phi_1$ \\
% F$(\phi)$ & $=$ & $true \, U \phi$  & \hspace{-0.4cm} in at least one \textbf{f}uture state \\
% G$(\phi)$ & $=$ & $\neg F \neg \phi$  & \hspace{-0.4cm} \textbf{g}lobally in all future states \\
%\end{tabular}

\subsection{Metric Temporal Logic}

\subsection{Signal Temporal Logic}